\section{Método} \label{sec:metodo}

\subsection{Diretrizes, PICOC e Questões de Pesquisa}

Conduzimos um mapeamento rápido descritivo (\textit{rapid mapping study}),
combinando o protocolo de~\cite{fernandes_interoperability_2018}
(GQM, PICOC, critérios e \textit{snowballing}) às adaptações rápidas
de~\cite{garritty_cochrane_2021}, com apoio da ferramenta Parsifal. Pelo PICOC,
adotamos \textit{Population} (\textit{Parkinson, Parkinsonism}),
\textit{Intervention} (\textit{PADS, smartwatch, wearable, inertial sensor,
accelerometer, gyroscope, ML/DL}), \textit{Comparison} (não se aplica, N/A),
\textit{Outcome} (\textit{classification, detection, diagnosis, differential
diagnosis, accuracy, F1-score}, AUC - área sob a curva ROC, \textit{sensitivity,
specificity}) e \textit{Context} (\textit{movement tasks, digital biomarkers, open dataset}). Sob o GQM, respondemos a \textbf{QP1}
(estudos, ano, veículo, contribuição), \textbf{QP2} (tarefas e recortes),
\textbf{QP3} (pré-processamento, atributos, modelos), \textbf{QP4} (validação,
métricas, desempenho, qualidade e reprodutibilidade) e \textbf{QP5} (lacunas e
oportunidades).

\subsection{Busca, Fontes e Fluxo de Seleção}

{\sloppy\setlength{\emergencystretch}{2em}
A busca abrangeu seis fontes: rastreamento de citações ao
artigo-base~\cite{varghese_machine_2024} no Google Scholar e na Nature Citations
(Dimensions), além de consultas por \textit{string} nas bases IEEE Xplore,
PubMed Central, ScienceDirect e Scopus. A \textit{string} geral combinou cinco
blocos de termos: \texttt{"PADS"}, \texttt{Parkinson*}, sensores
(\texttt{smartwatch}, \texttt{accelerometer}, \texttt{gyroscope},
\texttt{inertial sensor}, \texttt{wearable}), técnicas de IA (\texttt{accurac*},
\texttt{artificial intelligence}, \texttt{attention}, \texttt{framework},
\texttt{learning}, \texttt{model}, \texttt{network}) e tarefa
(\texttt{classification}, \texttt{detection}, \texttt{disease identification},
\texttt{screening}), ajustada a cada biblioteca (na ScienceDirect:
título, resumo e palavras-chave), em busca de
12/11/2026 consolidada até 13/11/2026.\par
}

Os \nImportados\ registros importados no Parsifal vieram de três rotas. Na (1)
\textbf{busca por citações}, obtivemos \nScholarDireto\ registros diretos no
Google Scholar e \nNatureExcl\ exclusivos da Nature Citations (\nNatureTotal\ no
total dessa fonte, sendo \nComumScholarNature\ em comum com o Scholar,
\nSoScholar\ exclusivos do Scholar e união de \nUniaoScholarNature). Nas (2)
\textbf{consultas por \textit{string}}, recuperamos \nStringsTotal\ registros
(IEEE \nStringsIEEE, PMC \nStringsPMC, SD \nStringsSD\ e Scopus \nStringsScopus).
No (3) \textbf{\textit{forward snowballing}} exclusivo (sem etapa
\textit{backward}, pois os usuários sucedem o PADS (2024), em adaptação
prospectiva de~\cite{fernandes_interoperability_2018}),
identificamos \nSnowballing\ registros únicos a partir de \nSementesSnowballing\
sementes. No Parsifal, esses \nSnowballing\ registros do \textit{snowballing}
somaram-se ao Google Scholar (\nScholarDireto\ + \nSnowballing\ =
\nScholarBalde), totalizando \nScholarBalde\ + \nNatureExcl\ + \nStringsTotal\ =
\nImportados\ importações.

Adotamos quatro critérios de inclusão (\textbf{CI1}: uso do PADS, \textbf{CI2}:
publicação entre o lançamento do PADS, em 05/01/2024 no artigo e 25/03/2024 no
PhysioNet, e 12/11/2026, \textbf{CI3}: artigo completo revisado por pares e
\textbf{CI4}: inglês ou português) e seis de exclusão (\textbf{CE1}: citação
apenas relacionada, \textbf{CE2}: literatura cinzenta (sem revisão formal: \textit{preprints}, livros e resumos), \textbf{CE3}:
revisões/\textit{surveys}, \textbf{CE4}: duplicatas ou versões preliminares, \textbf{CE5}: sem texto completo e \textbf{CE6}: outro
idioma). Dos \nImportados\ registros, removemos \nDuplicados\ duplicados e
triamos \nTriados, resultando em \nRejeitados\ rejeitados (\nRejCEUm\ por
CI1/CE1, \nRejCETres\ por CE3 e \nRejCEDois\ por CE2, sendo \nRejMetadados\ por
metadados e \nRejTexto\ em texto completo, 0 isoladas por CE5 ou CE6) e
\nAceitos\ incluídos (Figura~\ref{fig:fluxo}).

\begin{figure}[htbp]
\centering
\begin{tikzpicture}[
  node distance=0.18cm and 0.20cm,
  box/.style={rectangle, draw=black, line width=0.65pt, align=center, font=\scriptsize, minimum height=0.42cm, inner sep=1.8pt},
  widebox/.style={rectangle, draw=black, line width=0.65pt, align=center, font=\scriptsize, minimum height=0.38cm, text width=7.0cm, inner sep=1.8pt},
  sidebox/.style={rectangle, draw=black, dashed, line width=0.55pt, align=left, font=\scriptsize, minimum height=0.38cm, text width=4.1cm, inner sep=1.8pt},
  arr/.style={-{Stealth[length=2mm]}, line width=0.65pt}
]
\node[box, text width=4.20cm] (r1) {\textbf{Rota 1: Citações base}\\Scholar: \nScholarDireto, Nature: \nNatureExcl\ excl.\ (\nNatureTotal)};
\node[box, right=0.20cm of r1, text width=4.20cm] (r2) {\textbf{Rota 2: Strings (\nStringsTotal)}\\IEEE \nStringsIEEE, PMC \nStringsPMC, SD \nStringsSD, Scopus \nStringsScopus};
\node[box, right=0.20cm of r2, text width=3.70cm] (r3) {\textbf{Rota 3: Forward Snowballing}\\ \nSnowballing\ reg.\ (\nSementesSnowballing\ sementes)};

\node[widebox, below=0.22cm of r2] (imp) {\textbf{Importados no Parsifal ($n=\nImportados$)} (\nScholarDireto\ + \nNatureExcl\ + \nSnowballing\ + \nStringsTotal)};
\node[widebox, below=0.16cm of imp] (tri) {\textbf{Registros únicos triados ($n=\nTriados$)}};
\node[sidebox, right=0.35cm of tri] (dup) {\textbf{Duplicados removidos: \nDuplicados} (CE4)};
\node[widebox, below=0.32cm of tri] (ace) {\textbf{Incluídos na síntese ($n=\nAceitos$)}: \nRevisadosPares\ revisados + \nPreprints\ \textit{preprints}};
\node[sidebox, right=0.35cm of ace] (rej) {\textbf{Rejeitados: \nRejeitados}\\CE1/CI1: \nRejCEUm, CE3: \nRejCETres, CE2: \nRejCEDois\\(\nRejMetadados\ metadados, \nRejTexto\ texto compl.)};

\draw[arr] (r1.south) |- ([yshift=0.11cm]imp.north) -- (imp.north);
\draw[arr] (r2.south) -- (imp.north);
\draw[arr] (r3.south) |- ([yshift=0.11cm]imp.north) -- (imp.north);
\draw[arr] (imp.south) -- (tri.north);
\draw[arr] (tri.east) -- (dup.west);
\draw[arr] (tri.south) -- (ace.north);
\draw[arr] ([yshift=-0.16cm]tri.south) -| ([xshift=-0.17cm]rej.west) -- (rej.west);
\end{tikzpicture}
\caption{Fluxo de identificação, triagem e inclusão dos \nTotal\ estudos por rota.}
\label{fig:fluxo}
\end{figure}

\subsection{Desvio do Protocolo, Qualidade e Extração}\label{sec:desvio_qa}

\textbf{Desvio declarado e casos limítrofes:} apesar de CE2 e CI3 excluírem
literatura cinzenta, mantivemos os \nPreprints\ \textit{preprints} localizados como
desvio declarado do protocolo. Em uma base recente (2024) e com corpus revisado
ainda reduzido ($n=\nRevisadosPares$), esses manuscritos de 2026 atendem a CI1 e
CI4, alcançam qualidade metodológica mínima ($\ge 70\%$, faixa de 70\% a 90\%) e
registram debates sobre vazamento em \textit{augmentation}~\cite{djemaa_quantifying_2026},
sem versão publicada até o corte. Os outros dez registros rejeitados em CE2
permaneceram excluídos (nove sem uso confirmado do PADS em CI1 e um resumo de
pôster). Esses \nPreprints\ \textit{preprints},
incluindo~\cite{zannat_label-efficient_2026} (risco alto de vazamento por
janelamento, segmentação temporal das séries, dependente de classe), recebem
análise isolada e de sensibilidade na Seção~\ref{sec:preprints}. Mantivemos
também o artigo-semente~\cite{varghese_machine_2024} e dois estudos restritos ao
PDNMS~\cite{padilha_parkinsons_2024,ouabd_parkinsons_2026}, pois~\cite{padilha_parkinsons_2024}
é um resumo expandido revisado por pares de duas páginas com experimento e métricas
completas.

\textbf{Qualidade metodológica (QA1--QA5):} avaliamos definição da tarefa (QA1),
recorte (QA2), prevenção de vazamento entre sujeitos (QA3), métricas ao
desbalanceamento (QA4) e replicabilidade/código (QA5) em \textit{Yes}=1,0,
\textit{Partially}=0,5 e \textit{No}=0,0, como escore descritivo distinto do
risco de vazamento. Itens \textit{N/A} (\cite{abola_differentiating_2025} em
análise estatística e \cite{ren_baseline_2026} em linha de base sem modelo
preditivo, avaliado apenas em QA2 e QA5 com 0/2) saíram do escore proporcional
(\texttt{nota\_qa\_pct}), agrupado nas faixas Alta ($\ge 80\%$), Média
($50\%$--$79\%$) e Baixa ($<50\%$).

\textbf{Extração e risco de vazamento:} o autor extraiu os 38 campos no Parsifal
e realizou segunda conferência integral, sem dupla extração nem cálculo de
$\kappa$ (\cite{garritty_cochrane_2021}), validando tipo e veículo no Zotero e
nos PDFs. Segundo a literatura de
vazamento~\cite{kapoor_leakage_2023,vandewiele_overly_2021,djemaa_quantifying_2026},
o risco reflete o relato: \textit{N/A} (sem modelo preditivo), \textit{Alto}
(divisão por janela/gravação, mistura de sujeitos ou balanceamento de classes e
pré-processamento pré-\textit{split}, anteriores à divisão treino/teste),
\textit{Baixo} (divisão por sujeito e balanceamento/pré-processamento ausentes ou
restritos ao treino), \textit{Moderado} (parcial ou sem detalhe do
pré-processamento) e \textit{Indeterminado} (sem dados).
